\begin{abstract}

Do large language models (LLMs) internalize the partisan structure of American political discourse in a way that mirrors empirical patterns in public opinion? We investigate this question by comparing two quantities across 199 political topics from the General Social Survey (GSS, 2021--2024): (1) \emph{activation-space partisan divergence}, the unsupervised Mahalanobis distance between Democrat and Republican centroids in per-head PCA-reduced attention-activation space when models simulate political actors; and (2) \emph{survey partisanship}, the normalized partisan gap in GSS response means. Across 17 transformer models (3B--72B parameters) and three prompt conditions, we find consistent positive correlations between partisan divergence and survey partisanship, with the strongest alignment reaching $r = 0.74$ on 126 public-issue topics. Three findings emerge. First, larger models achieve higher alignment, with 70B+ models consistently surpassing $r = 0.65$. Second, a 24B uncensored fine-tune (Dolphin3.0-Mistral-24B) outperforms safety-aligned models with three times the parameter count; a controlled comparison at identical scale---the same base model with safety alignment partially restored (Dolphin-24B-Venice)---scores 0.11 points lower, directly isolating the RLHF suppression effect from parameter count. This suggests that RLHF attenuates but does not erase the partisan geometry encoded during pretraining. Third, politician-name prompts yield substantially higher alignment than demographic-persona prompts constructed from real respondent profiles, indicating that models encode political knowledge more strongly through named-entity associations than through demographic inference. The alignment extends to 73 private-life topics (religious practice, sexual attitudes, moral views), with some models matching or exceeding their public-issue scores, suggesting that LLMs have internalized not only issue partisanship but the broader cultural sorting of American society. These findings are descriptive and correlational.
